%% Example CV - Juan David Posso Rengifo
%% Comprehensive overview of all competencies, experience, and achievements.
%% Use as a master reference when tailoring role-specific CVs.
%%
%% Master populated via /setup; tailor and translate to the language of each posting.
%% Compile with: lualatex main_example.tex
%% (pdflatex often fails on modern MiKTeX with fontawesome5 font-expansion errors.)

\documentclass[11pt,a4paper,sans]{moderncv}
\moderncvstyle{banking}
\moderncvcolor{blue}

% Force the name and section headings to render in moderncv blue (color1).
% Default banking leaves them black: moderncvstylebanking.sty's \colorlet
% copies (not aliases) the pre-scheme accent colour, so the name colours are
% frozen before \moderncvcolor runs. Re-let them after. \namefont is the hook
% every name-style macro routes through, so this also works on moderncv 2.3.1
% (Debian/Ubuntu apt), which has no \firstnamestyle/\lastnamestyle at all.
\renewcommand*{\namefont}{\fontsize{25}{28}\bfseries\upshape}
\colorlet{firstnamecolor}{color1}
\colorlet{lastnamecolor}{color1}
\colorlet{namecolor}{color1}
\renewcommand*{\sectionstyle}[1]{{\sectionfont\color{color1}#1}}

\usepackage[utf8]{inputenc}
% pdflatex fallback only (the documented engine is lualatex, which skips this
% branch). Without T1 font encoding pdflatex builds accented letters with
% \accent, and the PDF text layer stores them decomposed - `e` + U+0300 rather
% than U+00E8 - so an ATS keyword match on "Genève" fails while the page looks
% right. moderncv 2.5 loads T1 itself under pdflatex; 2.3.1 (Debian/Ubuntu apt)
% does not. \ifpdftex comes from iftex, which every moderncv version loads.
\ifpdftex\usepackage[T1]{fontenc}\fi
% moderncv loads hyperref itself in an \AtEndPreamble hook, so \hypersetup
% must go in an \AtEndPreamble of our own: on moderncv < 2.4 a top-level
% \usepackage{hyperref} clashes with the class's own
% \RequirePackage[unicode]{hyperref}. From 2.4.0 the class passes its options
% through \PassOptionsToPackage instead, which is what removes that clash.
\AtEndPreamble{\hypersetup{
    colorlinks=true,
    linkcolor=blue,
    filecolor=magenta,
    urlcolor=blue,
    pdftitle={Juan David Posso Rengifo - CV},
    % Keep pdfpagemode=UseNone: this block runs after moderncv's own
    % \AtEndPreamble (moderncv.cls sets pdfpagemode there), so a FullScreen
    % value here would win and open every CV in fullscreen presentation mode.
    pdfpagemode=UseNone,
}}
\usepackage[scale=0.80]{geometry}
\usepackage{import}

% personal data
\name{Juan David}{Posso Rengifo}
% If you have no address to list, DELETE this whole line. \address{}{}{} fails
% with "There's no line here to end" on every moderncv version.
\address{Santiago de Cali, Colombia}{}{}
\phone[mobile]{+57 3145291839}
\email{juandapore3098@gmail.com}
\extrainfo{\href{https://www.linkedin.com/in/juandavidposso/}{LinkedIn}, \href{https://github.com/juanda3098}{GitHub}}

\usepackage{needspace}
\begin{document}
\makecvtitle

\section{Perfil}
Ingeniero Multimedia con experiencia en desarrollo frontend y backend y fortalezas en React, Apex y Aura Components. Participación en un portal de autoservicio en producción y en las etapas iniciales de una migración de Salesforce a un CRM propio. Experiencia liderando un equipo de tres desarrolladores.

\section{Competencias técnicas}
\cvitem{Frontend}{React, JavaScript, Redux, Lightning Web Components y Aura Components.}
\cvitem{Backend}{Apex, Node.js, NestJS y MySQL.}
\cvitem{Salesforce}{Salesforce Development, Visualforce y Salesforce Flow.}
\cvitem{Prácticas}{Domain-Driven Design, arquitectura hexagonal, revisión de código, Playwright y liderazgo técnico.}
\cvitem{Herramientas de IA}{Codex para desarrollo, revisión de código y creación de skills.}

\section{Experiencia profesional}
\needspace{5\baselineskip}
\cventry{enero 2025 - septiembre 2026}{Software Engineer III}{Veevart}{Cali, Colombia; híbrido}{}{
\begin{itemize}
\item Desarrollo de proyectos completos desde el análisis técnico hasta la implementación.
\item Levantamiento y modelado del dominio con DDD para apoyar el inicio de la transición de Salesforce a una plataforma CRM propia.
\item Uso de arquitectura hexagonal en el desarrollo de software.
\item Revisión de código y acompañamiento en la resolución de problemas técnicos.
\item Uso de Codex para desarrollo, revisiones y creación de skills orientadas a mejorar procesos del equipo.
\end{itemize}}
\needspace{5\baselineskip}
\cventry{junio 2023 - enero 2025}{Lead Engineer}{Veevart}{Cali, Colombia; híbrido}{}{
\begin{itemize}
\item Liderazgo de tres desarrolladores, en colaboración con un Product Manager y con apoyo de QA.
\item Acompañamiento técnico al equipo para resolver problemas durante el desarrollo de aplicaciones.
\end{itemize}}

\newpage
\section{Experiencia profesional (continuación)}
\needspace{5\baselineskip}
\cventry{enero 2023 - junio 2023}{Software Engineer II}{Veevart}{Cali, Colombia; híbrido}{}{}
\needspace{5\baselineskip}
\cventry{octubre 2021 - enero 2023}{Software Engineer I}{Veevart}{Cali, Colombia; híbrido}{}{}

\subsection{Proyecto destacado en Veevart}
\cvitem{Portal en producción}{Desarrollo frontend y backend de un portal donde los usuarios de los clientes gestionaban membresías, pagos recurrentes y donaciones, incluyendo el cambio de la tarjeta de crédito asociada a membresías o donaciones.}

\needspace{5\baselineskip}
\cventry{abril 2021 - septiembre 2021}{Desarrollador junior}{Kaleido Lab}{}{Fechas aproximadas}{
\begin{itemize}
\item Desarrollo de la estructura de experiencias web interactivas para espacios virtuales de reuniones y tiendas con compras online interactivas.
\item Uso de JavaScript, React y MySQL.
\end{itemize}}

\section{Educación}
\needspace{5\baselineskip}
\cventry{2015 - 2020}{Ingeniero Multimedia}{Universidad Autónoma de Occidente}{}{Graduado}{
Proyecto de grado: desarrollo de una plataforma para la gestión de pacientes de un centro de fisioterapia.}

\section{Idiomas}
\cvitem{Español}{Nativo.}
\cvitem{Inglés}{B1.}
\end{document}
